\section{Interface matching systems}
\label{app:interface-matching}

The local interface maps used in Sec.~\ref{sec:conductance} are obtained from selector matrices that extract the active layer and the terminated-layer components from the central state $\Xi$.  At the left contact, layer~1 continues into the ferromagnetic lead and layer~2 terminates.  Let $S_L$ extract the four layer-1 BdG components and $M_L$ extract the two independent missing-neighbor components of the terminated layer.  Then the homogeneous left-interface problem is
\begin{equation}
    \begin{pmatrix}
        S_LV_+ & -Z_F^-\\
        M_LV_+ & 0_{2\times2}
    \end{pmatrix}
    \begin{pmatrix}
        a_+^L\\ r
    \end{pmatrix}
    =
    -
    \begin{pmatrix}
        S_LV_-\\ M_LV_-
    \end{pmatrix}a_-^L.
    \label{eq:left-interface-system-app}
\end{equation}
Adding the physical incident state replaces the right-hand side of the active-layer equation by $\zeta_{\rm in}-S_LV_-a_-^L$, leading to the affine map in Eq.~\eqref{eq:left-interface-map}.

At the right contact, the active selector depends on the geometry.  For $g=\oneone$, the superconducting lead continues from layer~1; for $g=\onetwo$, it continues from layer~2.  The right-interface system is
\begin{equation}
    \begin{pmatrix}
        S_R^{(g)}V_- & -Z_S^+\\
        M_R^{(g)}V_- & 0_{2\times2}
    \end{pmatrix}
    \begin{pmatrix}
        a_-^R\\ t
    \end{pmatrix}
    =
    -
    \begin{pmatrix}
        S_R^{(g)}V_+\\ M_R^{(g)}V_+
    \end{pmatrix}a_+^R.
    \label{eq:right-interface-system-app}
\end{equation}
The full $6\times6$ systems in Eqs.~\eqref{eq:left-interface-system-app} and \eqref{eq:right-interface-system-app}, solved with pivoted linear algebra, are the numerical definition of the contact maps.  Schur-complement reductions are useful analytically, but a singular chosen pivot does not by itself imply that the full contact problem is singular.
